\begin{helper}[Minus-closedness transport for controlled affected colors]
Let $G$ be a finite loopless multigraph with palette $[k]$, distinguished
edge $e_0$, and exact tree-sequence $U$.  Let $\varphi$ be proper and
leave exactly $e_0$ uncolored.  Let $\alpha\ne\beta$ and
$\sigma=\noderef{TwoColorComponentSwap}(G,\varphi,\alpha,\beta,z)$.
Let $R\subseteq[k]$ and suppose
\noderef{TreeColorSetMinusClosed}$(G,\varphi,U,R)$ holds.
Suppose all vertex missing sets on $U$ agree under $\sigma,\varphi$.
Assume that for each $c\in\{\alpha,\beta\}$, if $c$ is missing on $U$
and singleton-closed there under $\varphi$, the recoloring is an
\noderef{OutsideKempeRecoloring} of $U$.
Then \noderef{TreeColorSetMinusClosed}$(G,\sigma,U,R)$ holds.
\end{helper}
\begin{proof}
The vertexwise equality gives
$\overline\sigma(U)=\overline\varphi(U)$ by
\noderef{TreeMissingColors}.
Fix $c\in\overline\sigma(U)-R$.  It belongs to
$\overline\varphi(U)-R$, so old minus-closedness says its old boundary
is empty, by \noderef{TreeColorSetMinusClosed} and
\noderef{TreeColorSetClosed}.  In particular $c$ is singleton-closed.
Suppose an exact boundary edge $f$ has new color $c$.
If $c$ is one of the affected colors, the stated outside-recoloring
assumption applies because its old missingness and singleton closedness
were just proved.  The endpoint characterization of
\noderef{BoundaryEdges} makes $f$ incident with $U$.
\noderef{OutsideKempeRecoloringStable} therefore fixes its exact value,
so its old color was also $c$.
If $c$ is unaffected, its edge-color membership is unchanged by
\noderef{TwoColorComponentSwap}: inside the component only $\alpha$
and $\beta$ values exchange, and outside it every value is fixed.
Thus its old color was $c$ in this case too.
Boundary membership depends only on endpoints.  We have found an old
$c$-colored boundary edge, contradicting old closedness.
Every color of $\overline\sigma(U)-R$ therefore has empty new boundary,
which is the required minus-closedness.  The argument does not require
that every affected color lie in $R$, or that the entire component avoid
$U$ without the explicit closed-color premise.
\end{proof}

